\Needspace{7\baselineskip}\SourceLeaf{85}

\noindent
Eliminiert man daraus $\dfrac{dX}{dx}$, so ergibt sich

\[
\frac{dx}{du} = X\frac{dx}{dt}.
\]

\noindent
Da nun

\[
\frac{d\varphi x}{dt} = \frac{d\varphi x}{dx}\frac{dx}{dt} \qquad \text{und} \qquad
\frac{d\varphi x}{du} = \frac{d\varphi x}{dx}\frac{dx}{du},
\]

\noindent
gilt:

\begin{equation}\tag{1}
\frac{d\varphi x}{du} = X\frac{d\varphi x}{dt},
\end{equation}

\noindent
folglich auch, wenn man $\varphi x$ durch $X$ ersetzt (was zulässig ist),

\begin{equation}\tag{2}
\frac{dX}{du} = X\frac{dX}{dt},
\end{equation}

\Needspace{12\baselineskip}\noindent
Differenziert man (1) erneut nach $t$ und $u$, so erhält man

\begin{equation}\tag{3}
\frac{d^2\varphi x}{dt\,du}=\frac{dX}{dt}\frac{d\varphi x}{dt}+X\frac{d^2\varphi x}{dt^2},
\end{equation}
\[
\frac{d^2\varphi x}{du^2}=\frac{dX}{du}\frac{d\varphi x}{dt}+X\frac{d^2\varphi x}{dt\,du},
\]

\noindent
Setzt man hier die gefundenen Werte von $\dfrac{d^2\varphi x}{dt\,du}$ und $\dfrac{dX}{du}$
ein, so ergibt sich

\[
\frac{d^2\varphi x}{du^2} = 2X\frac{dX}{dt}\frac{d\varphi x}{dt} + X^2\frac{d^2\varphi x}{dt^2}
\]

\noindent
oder

\begin{equation}\tag{4}
\frac{d^2\varphi x}{du^2} = \frac{dX^2 \frac{d\varphi x}{dt}}{dt}.
\end{equation}

\noindent
Differenziert man diese Gleichung wiederum nach $u$, so ergibt sich

\[
\frac{d^3\varphi x}{du^3} = \frac{d^2X^2 \frac{d\varphi x}{dt}}{dt\,du}.
\]

\noindent\hspace*{\parindent}Da es aber bekanntlich gleichgültig ist, ob man $X^2\dfrac{d\varphi x}{dt}$ erst nach
$t$ und dann nach $u$ differenziert oder in umgekehrter Reihenfolge, und

\[
\frac{dX^2 \frac{d\varphi x}{dt}}{du}\mathord{\cdot} = 2X\frac{dX}{du}\frac{d\varphi x}{dt} + X^2\frac{d^2\varphi x}{dt\,du}
\]

\noindent
oder, wenn man die Werte von $\dfrac{dX}{du}$, $\dfrac{d^2\varphi x}{dt\,du}$ aus
(2) und (3) einsetzt,
